\documentclass[11pt,a4paper]{article}
\usepackage{fontspec,amsmath,amssymb,amsthm,microtype}
\setmainfont{lmroman10-regular.otf}[BoldFont=lmroman10-bold.otf,ItalicFont=lmroman10-italic.otf,BoldItalicFont=lmroman10-bolditalic.otf]
\usepackage{xcolor}
\usepackage[most]{tcolorbox}
\usepackage{mathrsfs}
\providecommand{\lt}{<}
\providecommand{\gt}{>}
\setmonofont{DejaVuSansMono.ttf}[BoldFont=DejaVuSansMono-Bold.ttf,ItalicFont=DejaVuSansMono-Oblique.ttf,BoldItalicFont=DejaVuSansMono-BoldOblique.ttf,Scale=0.86]
\newfontfamily\leanfallback{STIXTwoMath-Regular.otf}[Scale=0.86]
\definecolor{leanrule}{gray}{0.62}
\definecolor{leanback}{gray}{0.962}
\definecolor{leantitle}{gray}{0.22}
\definecolor{leanmuted}{gray}{0.38}
\newtcolorbox{leanbox}[1]{enhanced,breakable,colback=leanback,colframe=leanrule,
  boxrule=0.4pt,arc=0.6mm,left=2.4mm,right=2.4mm,top=1.8mm,bottom=1.8mm,
  before skip=0.8em,after skip=1em,fonttitle=\small\bfseries,coltitle=leantitle,
  colbacktitle=leanback,titlerule=0.3pt,title={#1}}
\newlength{\leancharwidth}
\newenvironment{leancode}{\par\smallskip\begingroup\ttfamily\footnotesize\raggedright
  \settowidth{\leancharwidth}{M}\setlength{\parskip}{0pt}\setlength{\parindent}{0pt}
  \frenchspacing}{\par\endgroup\smallskip}
\newcommand{\leanline}[2]{\hangindent=\dimexpr#1\leancharwidth+4\leancharwidth\relax
  \hangafter=1\noindent\hspace*{#1\leancharwidth}#2\par}
\newcommand{\leansource}[1]{\par{\footnotesize\color{leanmuted}#1}\par}
\newcommand{\leanlabel}[1]{\par\smallskip\noindent{\small\bfseries #1}\par}
\usepackage[margin=25mm]{geometry}
\usepackage[unicode,hidelinks]{hyperref}
\hypersetup{pdfauthor={LangevinLean Project}}
\setlength{\emergencystretch}{3em}
\setlength{\parskip}{0.4em}
\setlength{\parindent}{0pt}
\allowdisplaybreaks[2]
\raggedbottom
\title{\texorpdfstring{Strong convergence of bounded observables without stationary approximants}{Strong convergence of bounded observables without stationary approximants}}
\author{LangevinLean Project}
\date{}
\begin{document}\maketitle
Draft mathematical exposition; these notes may contain mistakes.


Euler approximation states need not have the target law. A bounded continuous observable nevertheless gives a genuine square-integrable random variable at each state, and convergence in measure of those states implies strong convergence of the bounded observable compositions.


\section{\texorpdfstring{Actual composition, subsequences and dominated convergence}{Actual composition, subsequences and dominated convergence}}


\textbf{Claim.} For a bounded continuous \(\phi\), if \(X_i\to X\) in measure on a probability space, then \(\phi(X_i)\to\phi(X)\) in \(L^2\), even when their marginal laws differ.



\textbf{Proof.} Boundedness gives square integrability and an exact representative for every composition. Every sequence approaching the filter has a subsequence along which state convergence holds almost surely. Continuity gives almost-sure convergence of the compositions. Their squared difference is bounded by the integrable constant \(4\|\phi\|_\infty^2\), so dominated convergence forces its integral to zero along that subsequence. The subsequence criterion gives the full strong limit. Strong state convergence is a special case because it implies convergence in measure.

\begin{leanbox}{Lean: \texttt{boundedComposition}, \texttt{boundedComposition\_\allowbreak{}coe}, \texttt{boundedComposition\_\allowbreak{}sub\_\allowbreak{}norm\_\allowbreak{}sq} and 2 more}
\leansource{\texttt{boundedComposition} \quad WassersteinMixing/\allowbreak{}BoundedComposition.lean, lines 23–25}
\begin{leancode}
\leanline{0}{def\ boundedComposition\ (X\ :\ Lp\ E\ 2\ Q)\ (φ\ :\ E\ →ᵇ\ H)\ :\ Lp\ H\ 2\ Q\ :=}
\leanline{2}{(MemLp.of\_bound\ (φ.continuous.comp\_aestronglyMeasurable\ (Lp.aestronglyMeasurable\ X))}
\leanline{4}{‖φ‖\ (Eventually.of\_forall\ fun\ ω\ ↦\ φ.norm\_coe\_le\_norm\ (X\ ω))).toLp\ (fun\ ω\ ↦\ φ\ (X\ ω))}
\end{leancode}
\leansource{\texttt{boundedComposition\_\allowbreak{}coe} \quad WassersteinMixing/\allowbreak{}BoundedComposition.lean, lines 27–31}
\begin{leancode}
\leanline{0}{theorem\ boundedComposition\_coe\ (X\ :\ Lp\ E\ 2\ Q)\ (φ\ :\ E\ →ᵇ\ H)\ :}
\leanline{4}{boundedComposition\ Q\ X\ φ\ =ᵐ[Q]\ fun\ ω\ ↦\ φ\ (X\ ω)}
\end{leancode}
\leansource{\texttt{boundedComposition\_\allowbreak{}sub\_\allowbreak{}norm\_\allowbreak{}sq} \quad WassersteinMixing/\allowbreak{}BoundedComposition.lean, lines 32–41}
\begin{leancode}
\leanline{0}{private\ theorem\ boundedComposition\_sub\_norm\_sq\ (X\ Y\ :\ Lp\ E\ 2\ Q)\ (φ\ :\ E\ →ᵇ\ H)\ :}
\leanline{4}{‖boundedComposition\ Q\ X\ φ\ -\ boundedComposition\ Q\ Y\ φ‖\ \textasciicircum{}\ 2\ =}
\leanline{6}{∫\ ω,\ ‖φ\ (X\ ω)\ -\ φ\ (Y\ ω)‖\ \textasciicircum{}\ 2\ ∂Q}
\end{leancode}
\leansource{\texttt{boundedComposition\_\allowbreak{}tendsto\_\allowbreak{}inMeasure} \quad WassersteinMixing/\allowbreak{}BoundedComposition.lean, lines 44–79}
\begin{leancode}
\leanline{0}{theorem\ boundedComposition\_tendsto\_inMeasure\ \{I\ :\ Type*\}\ \{l\ :\ Filter\ I\}}
\leanline{4}{[l.IsCountablyGenerated]\ (X\ :\ I\ →\ Lp\ E\ 2\ Q)\ (Y\ :\ Lp\ E\ 2\ Q)}
\leanline{4}{(hlim\ :\ TendstoInMeasure\ Q\ (fun\ i\ ω\ ↦\ X\ i\ ω)\ l\ Y)\ (φ\ :\ E\ →ᵇ\ H)\ :}
\leanline{4}{Tendsto\ (fun\ i\ ↦\ boundedComposition\ Q\ (X\ i)\ φ)\ l\ ({\leanfallback 𝓝}\ (boundedComposition\ Q\ Y\ φ))}
\end{leancode}
\leansource{\texttt{boundedComposition\_\allowbreak{}tendsto} \quad WassersteinMixing/\allowbreak{}BoundedComposition.lean, lines 81–84}
\begin{leancode}
\leanline{0}{theorem\ boundedComposition\_tendsto\ \{I\ :\ Type*\}\ \{l\ :\ Filter\ I\}\ [l.IsCountablyGenerated]}
\leanline{4}{(X\ :\ I\ →\ Lp\ E\ 2\ Q)\ (Y\ :\ Lp\ E\ 2\ Q)\ (hlim\ :\ Tendsto\ X\ l\ ({\leanfallback 𝓝}\ Y))\ (φ\ :\ E\ →ᵇ\ H)\ :}
\leanline{4}{Tendsto\ (fun\ i\ ↦\ boundedComposition\ Q\ (X\ i)\ φ)\ l\ ({\leanfallback 𝓝}\ (boundedComposition\ Q\ Y\ φ))}
\end{leancode}
\leanlabel{How the Lean proof follows the argument}
\texttt{bounded\allowbreak{}Composition} constructs the actual \(L^2\) class from its uniform bound. The private norm identity uses its exact representative. \texttt{bounded\allowbreak{}Composition\_\allowbreak{}tendsto\_\allowbreak{}in\allowbreak{}Measure} applies subsequence extraction and dominated convergence; the final theorem converts strong state convergence to convergence in measure.
\end{leanbox}

\end{document}
